\chapter{ФРАГМЕНТИ КОДУ}

\section{TimescaleDB Schema}

\lstset{style=code}
\begin{lstlisting}[language=SQL, caption={Повна схема бази даних}]
-- Companies (tenants)
CREATE TABLE companies (
    id UUID PRIMARY KEY DEFAULT gen_random_uuid(),
    name VARCHAR(255) NOT NULL UNIQUE,
    created_at TIMESTAMPTZ DEFAULT NOW()
);

-- Equipment
CREATE TABLE equipment (
    id UUID PRIMARY KEY DEFAULT gen_random_uuid(),
    company_id UUID REFERENCES companies(id),
    name VARCHAR(255) NOT NULL,
    type VARCHAR(100) NOT NULL,
    location VARCHAR(255),
    created_at TIMESTAMPTZ DEFAULT NOW()
);

-- Sensors
CREATE TABLE sensors (
    id UUID PRIMARY KEY DEFAULT gen_random_uuid(),
    equipment_id UUID REFERENCES equipment(id),
    company_id UUID REFERENCES companies(id),
    name VARCHAR(255) NOT NULL,
    sensor_type VARCHAR(50) NOT NULL,
    unit VARCHAR(50),
    created_at TIMESTAMPTZ DEFAULT NOW()
);

-- Sensor Measurements (Hypertable)
CREATE TABLE sensor_measurements (
    timestamp TIMESTAMPTZ NOT NULL,
    sensor_id UUID NOT NULL REFERENCES sensors(id),
    equipment_id UUID NOT NULL REFERENCES equipment(id),
    company_id UUID NOT NULL REFERENCES companies(id),
    value DOUBLE PRECISION NOT NULL,
    quality_code INTEGER DEFAULT 192
);

SELECT create_hypertable('sensor_measurements', 'timestamp');

-- Continuous Aggregates
CREATE MATERIALIZED VIEW sensor_agg_1min
WITH (timescaledb.continuous) AS
SELECT 
    time_bucket('1 minute', timestamp) AS bucket,
    sensor_id,
    equipment_id,
    company_id,
    AVG(value) AS avg_value,
    MIN(value) AS min_value,
    MAX(value) AS max_value,
    COUNT(*) AS count
FROM sensor_measurements
GROUP BY bucket, sensor_id, equipment_id, company_id;
\end{lstlisting}

\section{Spark Streaming Job}

\begin{lstlisting}[language=Python, caption={Spark job для обробки raw даних}]
from pyspark.sql import SparkSession
from pyspark.sql.functions import from_json, col
from pyspark.sql.types import StructType, StringType, DoubleType

# Schema for incoming messages
schema = StructType() \
    .add("sensor_id", StringType()) \
    .add("equipment_id", StringType()) \
    .add("company_id", StringType()) \
    .add("timestamp", StringType()) \
    .add("value", DoubleType())

# Read from Kafka
df = spark \
    .readStream \
    .format("kafka") \
    .option("kafka.bootstrap.servers", "kafka:9092") \
    .option("subscribe", "sensor-data-raw") \
    .load()

# Parse JSON
parsed_df = df \
    .select(from_json(col("value").cast("string"), schema).alias("data")) \
    .select("data.*")

# Write to TimescaleDB
def write_to_db(batch_df, batch_id):
    batch_df.write \
        .format("jdbc") \
        .option("url", jdbc_url) \
        .option("dbtable", "sensor_measurements") \
        .option("user", db_user) \
        .option("password", db_password) \
        .mode("append") \
        .save()

query = parsed_df.writeStream \
    .foreachBatch(write_to_db) \
    .start()

query.awaitTermination()
\end{lstlisting}

\section{ML Model Training}

\begin{lstlisting}[language=Python, caption={Навчання Isolation Forest моделі}]
import pandas as pd
import numpy as np
from sklearn.ensemble import IsolationForest
from sklearn.preprocessing import StandardScaler
import mlflow

# Extract data
async def extract_training_data(company_id: str):
    query = """
    SELECT 
        timestamp,
        sensor_id,
        value,
        LAG(value, 1) OVER (PARTITION BY sensor_id ORDER BY timestamp) as prev_value,
        LAG(value, 2) OVER (PARTITION BY sensor_id ORDER BY timestamp) as prev_value_2
    FROM sensor_measurements
    WHERE company_id = %s
    ORDER BY timestamp
    """
    df = await db.fetch_all(query, company_id)
    return pd.DataFrame(df)

# Feature engineering
def create_features(df):
    df['value_change'] = df['value'] - df['prev_value']
    df['value_change_rate'] = df['value_change'] / df['prev_value']
    df['rolling_mean_5'] = df.groupby('sensor_id')['value'].transform(
        lambda x: x.rolling(5).mean()
    )
    df['rolling_std_5'] = df.groupby('sensor_id')['value'].transform(
        lambda x: x.rolling(5).std()
    )
    return df.dropna()

# Train model
def train_model(X_train, contamination=0.1):
    with mlflow.start_run():
        model = IsolationForest(
            n_estimators=100,
            contamination=contamination,
            random_state=42
        )
        
        scaler = StandardScaler()
        X_scaled = scaler.fit_transform(X_train)
        
        model.fit(X_scaled)
        
        mlflow.log_params({
            "n_estimators": 100,
            "contamination": contamination
        })
        
        mlflow.sklearn.log_model(model, "model")
        mlflow.sklearn.log_model(scaler, "scaler")
        
    return model, scaler
\end{lstlisting}

\section{FastAPI Endpoints}

\begin{lstlisting}[language=Python, caption={API Gateway endpoints}]
from fastapi import FastAPI, WebSocket, Depends
from fastapi.security import HTTPBearer, HTTPAuthorizationCredentials
import jwt

app = FastAPI()
security = HTTPBearer()

# JWT verification
def verify_token(credentials: HTTPAuthorizationCredentials):
    token = credentials.credentials
    payload = jwt.decode(token, SECRET_KEY, algorithms=["HS256"])
    return payload

# Get latest sensor data
@app.get("/api/sensors/{sensor_id}/latest")
async def get_latest_data(
    sensor_id: str,
    credentials: HTTPAuthorizationCredentials = Depends(security)
):
    user = verify_token(credentials)
    
    # Get from cache first
    cached = await redis.get(f"sensor:{sensor_id}")
    if cached:
        return json.loads(cached)
    
    # Fallback to database
    query = """
    SELECT * FROM sensor_measurements
    WHERE sensor_id = %s
    ORDER BY timestamp DESC
    LIMIT 1
    """
    result = await db.fetch_one(query, sensor_id)
    return result

# WebSocket for real-time data
@app.websocket("/ws/sensors")
async def websocket_endpoint(websocket: WebSocket):
    await websocket.accept()
    
    while True:
        # Get latest data from Redis
        data = await redis.xread({"sensor-stream": "$"})
        await websocket.send_json(data)
        await asyncio.sleep(1)
\end{lstlisting}
